\newcommand{\usecase}[6]{
\subsection*{#1}
\textbf{Actors principals:} #2\\
\textbf{Objectiu:} #3\\
\textbf{Precondicions:} #4\\
\textbf{Descripció:} #5\\
\textbf{Postcondicions:} #6
}

\documentclass[../main.tex]{subfiles}

\begin{document}


\section{Casos d’ús dels enquestats}
A continuació es descriuen els principals casos d’ús associats als usuaris enquestats, 
tant registrats com no registrats.


%%%%%%%%%%%%%%%%%%%%%%%%%%%%%%%%%%%%%%%%%%%%%%%%%%%%%%%%%%%%%%%%%%%%%%%%%%%%%%%%%%%%%%%%%%%
\subsection*{Cas d’ús 1: Respondre una enquesta}

\textbf{Actors principals:} Enquestat (registrat o no registrat)

\textbf{Objectiu:} Permetre a l’usuari respondre una enquesta disponible a la plataforma.

\textbf{Precondicions:}  
L’enquesta ha d’estar activa.
%%L’enquesta ha d’estar activa i accessible (pública o compartida amb l’usuari).

\textbf{Descripció:}
\begin{enumerate}
    \item L’usuari accedeix a la pàgina de l’enquesta.
    \item Llegeix les preguntes i introdueix les seves respostes.
    \item Confirma i envia les respostes.
\end{enumerate}

\textbf{Postcondicions:}  
Les respostes queden emmagatzemades al sistema i associades a l’enquesta.

\textbf{Excepcions:}
\begin{itemize}
    \item Error de validació de dades (respostes no desades).
    \item L’enquesta ja està tancada o no és accessible.
\end{itemize}


%%%%%%%%%%%%%%%%%%%%%%%%%%%%%%%%%%%%%%%%%%%%%%%%%%%%%%%%%%%%%%%%%%%%%%%%%%%%%%%%%%%%%%%%%%%
\subsection*{Cas d’ús 2: Registrar-se al sistema}

\textbf{Actors principals:} Enquestat (no registrat)

\textbf{Objectiu:} Permetre a un usuari no registrat crear un compte per accedir a funcionalitats addicionals (crear enquestes, rebre recomanacions, etc.).

\textbf{Precondicions:}  
L’usuari no ha de tenir un compte actiu al sistema.

\textbf{Descripció:}
\begin{enumerate}
    \item L’usuari accedeix a la secció de registre.
    \item Introdueix les seves dades (nom d’usuari, contrasenya, correu electrònic...).
    \item El sistema valida la informació i crea el compte.
\end{enumerate}

\textbf{Postcondicions:}  
El compte queda registrat i l’usuari pot iniciar sessió al sistema.

\textbf{Excepcions:}
\begin{itemize}
    \item El nom d’usuari ja està en ús.
    \item Error en la validació o manca d’algun camp obligatori.
\end{itemize}

%%%%%%%%%%%%%%%%%%%%%%%%%%%%%%%%%%%%%%%%%%%%%%%%%%%%%%%%%%%%%%%%%%%%%%%%%%%%%%%%%%%%%%%%%%%
\subsection*{Cas d’ús 3: Consultar resultats públics}

\textbf{Actors principals:} Enquestat (registrat o no registrat)

\textbf{Objectiu:} Permetre a l’usuari visualitzar els resultats de les enquestes marcades com a públiques.

\textbf{Precondicions:}  
L’enquesta ha de tenir els resultats configurats com a públics per part de l’administrador.

\textbf{Descripció:}
\begin{enumerate}
    \item L’usuari accedeix a la secció de resultats públics.
    \item Selecciona una enquesta.
    \item El sistema mostra un resum estadístic de les respostes.
\end{enumerate}

\textbf{Postcondicions:}  
L’usuari visualitza les dades agregades de l’enquesta seleccionada.

\textbf{Excepcions:}
\begin{itemize}
    \item No hi ha resultats disponibles o l’enquesta no és pública.
\end{itemize}


%%%%%%%%%%%%%%%%%%%%%%%%%%%%%%%%%%%%%%%%%%%%%%%%%%%%%%%%%%%%%%%%%%%%%%%%%%%%%%%%%%%%%%%%%%%
\subsection*{Cas d’ús 4: Rebre recomanacions d’enquestes}

\textbf{Actors principals:} Enquestat registrat

\textbf{Objectiu:} Permetre al sistema suggerir enquestes a l’usuari segons els seus interessos o respostes prèvies.

\textbf{Precondicions:}  
L’usuari ha d’estar registrat i haver respost almenys una enquesta anteriorment.

\textbf{Descripció:}
\begin{enumerate}
    \item L’usuari inicia sessió i accedeix al seu panell personal.
    \item El sistema analitza el seu historial de respostes.
    \item Es mostren recomanacions d’enquestes rellevants per l'usuari.
\end{enumerate}

\textbf{Postcondicions:}  
L’usuari pot seleccionar una de les enquestes recomanades i respondre-la.

\textbf{Excepcions:}
\begin{itemize}
    \item L’usuari no té activitat prèvia suficient per generar recomanacions.
\end{itemize}





\section{Casos d’ús dels administradors d’enquestes}

A continuació es descriuen els principals casos d’ús associats als administradors d’enquestes,
és a dir, usuaris registrats que han creat una o més enquestes i poden gestionar-ne el contingut i els resultats.


\subsection*{Cas d’ús 5: Crear una enquesta}

\textbf{Actors principals:} Administrador (usuari registrat)

\textbf{Objectiu:} Permetre a l’usuari crear una nova enquesta a la plataforma.

\textbf{Precondicions:}  
L’usuari ha d’estar autenticat al sistema.

\textbf{Descripció:}
\begin{enumerate}
    \item L’usuari accedeix a la secció de creació d’enquestes.
    \item Introdueix la informació bàsica (títol, descripció, visibilitat, etc.).
    \item Afegeix preguntes i possibles respostes.
    \item Desa i publica l’enquesta.
\end{enumerate}

\textbf{Postcondicions:}  
L’enquesta queda registrada al sistema i disponible segons la seva configuració (pública o privada).

\textbf{Excepcions:}
\begin{itemize}
    \item Falten camps obligatoris (títol, preguntes...).
\end{itemize}

\subsection*{Cas d’ús 6: Editar una enquesta}

\textbf{Actors principals:} Administrador (propietari de l’enquesta)

\textbf{Objectiu:} Permetre modificar el contingut o la configuració d’una enquesta existent.

\textbf{Precondicions:}  
L’enquesta ha de pertànyer a l’usuari i no pot estar tancada (si s’han rebut respostes, algunes modificacions poden estar restringides).

\textbf{Descripció:}
\begin{enumerate}
    \item L’usuari accedeix a la seva llista d’enquestes.
    \item Selecciona una enquesta per editar.
    \item Modifica preguntes, opcions o configuracions generals.
    \item Desa els canvis.
\end{enumerate}

\textbf{Postcondicions:}  
L’enquesta queda actualitzada segons les modificacions fetes.

\textbf{Excepcions:}
\begin{itemize}
    \item L’enquesta no existeix o no pertany a l’usuari.
    \item No es poden modificar preguntes després de rebre respostes.
\end{itemize}

\subsection*{Cas d’ús 7: Eliminar una enquesta}

\textbf{Actors principals:} Administrador (propietari de l’enquesta)

\textbf{Objectiu:} Permetre a l’usuari eliminar una enquesta pròpia i totes les seves respostes associades.

\textbf{Precondicions:}  
L’enquesta ha de pertànyer a l’usuari.

\textbf{Descripció:}
\begin{enumerate}
    \item L’usuari accedeix a la seva llista d’enquestes.
    \item Selecciona l’enquesta que vol eliminar.
    \item Confirma la seva eliminació.
\end{enumerate}

\textbf{Postcondicions:}  
L’enquesta i les seves respostes s’eliminen del sistema.

\textbf{Excepcions:}
\begin{itemize}
    \item L’enquesta no existeix o no pertany a l’usuari.
\end{itemize}


\subsection*{Cas d’ús 8: Consultar resultats de les seves enquestes}

\textbf{Actors principals:} Administrador (propietari de l’enquesta)

\textbf{Objectiu:} Permetre a l’usuari visualitzar els resultats i estadístiques de les seves enquestes.

\textbf{Precondicions:}  
L’enquesta ha de tenir respostes registrades.

\textbf{Descripció:}
\begin{enumerate}
    \item L’usuari accedeix a la seva enquesta.
    \item El sistema mostra un resum estadístic de les respostes rebudes.
    \item L’usuari pot exportar o filtrar els resultats.
\end{enumerate}

\textbf{Postcondicions:}  
L’usuari visualitza o descarrega els resultats de la seva enquesta.

\textbf{Excepcions:}
\begin{itemize}
    \item L’enquesta encara no té respostes.
    \item Error en la generació de les estadístiques.
\end{itemize}

\subsection*{Cas d’ús 9: Gestionar la visibilitat de l’enquesta}

\textbf{Actors principals:} Administrador (propietari de l’enquesta)

\textbf{Objectiu:} Permetre a l’usuari configurar si la seva enquesta és pública, privada o restringida a certs usuaris o enquestadors.

\textbf{Precondicions:}  
L’usuari ha de ser el propietari de l’enquesta.

\textbf{Descripció:}
\begin{enumerate}
    \item L’usuari accedeix a la configuració de l’enquesta.
    \item Selecciona el nivell de visibilitat desitjat.
    \item Desa la configuració.
\end{enumerate}

\textbf{Postcondicions:}  
La visibilitat de l’enquesta queda actualitzada segons la nova configuració.

\textbf{Excepcions:}
\begin{itemize}
    \item L’enquesta no és editable (per exemple, està sota revisió o moderació).
\end{itemize}

\section{Casos d’ús dels enquestadors}

Els enquestadors són usuaris designats per un o més administradors per ajudar en la difusió i recollida de respostes d’enquestes.
Poden administrar enquestes a altres persones presencialment, però no poden modificar el contingut de les preguntes.

\subsection*{Cas d’ús 10: Afegir respostes a una enquesta assignada}

\textbf{Actors principals:} Enquestador

\textbf{Objectiu:} Permetre a l’enquestador recollir respostes d’una enquesta en nom d’altres participants.

\textbf{Precondicions:}  
L’enquestador ha de tenir assignada l’enquesta per un administrador.

\textbf{Descripció:}
\begin{enumerate}
    \item L’enquestador accedeix a la llista d’enquestes assignades.
    \item Selecciona l’enquesta que vol administrar.
    \item Introdueix respostes en nom dels participants o comparteix l’enllaç per a respondre.
\end{enumerate}

\textbf{Postcondicions:}  
Les respostes queden registrades correctament al sistema, associades a l’enquestador o al participant corresponent.

\textbf{Excepcions:}
\begin{itemize}
    \item L’enquesta no està disponible o ha estat tancada per l’administrador.
\end{itemize}

\section{Casos d’ús dels moderadors}

Els moderadors tenen permisos globals dins del sistema.
S’encarreguen de vetllar pel compliment de les normes, revisar el contingut de les enquestes i gestionar usuaris quan sigui necessari.

\subsection*{Cas d’ús 11: Revisar una enquesta}

\textbf{Actors principals:} Moderador

\textbf{Objectiu:} Permetre al moderador revisar una enquesta creada per un administrador per assegurar que compleix les polítiques del sistema.

\textbf{Precondicions:}  
L’enquesta ha d'estar marcada com a denunciada.

\textbf{Descripció:}
\begin{enumerate}
    \item El moderador accedeix a la llista d’enquestes denunciades.
    \item Consulta el contingut i revisa les preguntes.
\end{enumerate}

\textbf{Postcondicions:}  
L’estat de l’enquesta s’actualitza (denuncia aprovada o rebutjada).

\textbf{Excepcions:}
\begin{itemize}
    \item Error en l’accés a la informació de l’enquesta.
\end{itemize}
\textbf{Relacions:}  
Aquest cas d'ús pot derivar en el \textit{Cas d'ús 13: Eliminar una enquesta} si el moderador detecta contingut inadequat.


\subsection*{Cas d’ús 12: Suspendre o eliminar una enquesta}

\textbf{Actors principals:} Moderador

\textbf{Objectiu:} Permetre al moderador eliminar una enquesta que incompleixi les normes de la plataforma.

\textbf{Precondicions:}  
El moderador ha d’haver revisat l’enquesta o tenir una denúncia activa sobre ella.

\textbf{Descripció:}
\begin{enumerate}
    \item El moderador selecciona l’enquesta en qüestió.
    \item Elimina definitivament l’enquesta.
    \item Confirma l’acció.
\end{enumerate}

\textbf{Postcondicions:}  
L’enquesta deixa d’estar disponible.

\textbf{Excepcions:}
\begin{itemize}
    \item L’enquesta ja ha estat eliminada o bloquejada per un altre moderador.
\end{itemize}

\subsection*{Cas d’ús 13: Gestionar usuaris}

\textbf{Actors principals:} Moderador

\textbf{Objectiu:} Permetre al moderador gestionar el comportament dels usuaris dins la plataforma.

\textbf{Precondicions:}  
L’usuari a gestionar ha d’estar registrat al sistema.

\textbf{Descripció:}
\begin{enumerate}
    \item El moderador accedeix al panell d’usuaris.
    \item Consulta el perfil o historial d’un usuari.
    \item Pot bloquejar l'usuari seleccionat.
\end{enumerate}

\textbf{Postcondicions:}  
Les mesures aplicades afecten l’estat de l’usuari.

\textbf{Excepcions:}
\begin{itemize}
    \item Error en l’accés al perfil d’usuari.
\end{itemize}




\end{document}
